\documentclass[12pt]{article}
\usepackage[utf8]{inputenc}
\usepackage[T1]{fontenc}
\usepackage{hyperref}
\title{Tangga Language Resources for K-12 Curriculum Development}
\author{Compiled by Claude Sonnet 5 (claude-sonnet-5) for Andr\'as Kornai}
\date{}
\begin{document}
\maketitle

\section{Introduction}
Tangga (ISO 639-3 tgg, Glottocode tang1348) is an Oceanic Austronesian language of New Ireland Province, Papua New Guinea, spoken on the Tanga and Feni islands and a few mainland villages. Speaker estimates range from about 5,800 to 11,000. It must not be confused with Tanggu (tgu), a non-Austronesian Ramu language of Madang Province. Varieties are highly divergent and are separately coded (Fanamaket bjp, Warwar Feni hrw).

\section{Primary Educational Resources}
None found.

\section{Linguistic and Reference Materials}
\begin{itemize}
\item Tangga language (Wikipedia, ISO 639 page): \url{https://wikipedia.com/wiki/ISO_639:tgg}
\item Tanga people (Wikipedia): \url{https://en.wikipedia.org/wiki/Tanga_people}
\item WALS entry: \url{https://wals.info/languoid/lect/wals_code_tga}
\item Cambridge Language Centre page on Tanga: \url{https://www.langcen.cam.ac.uk/resources/langt/tanga.html}
\item Verbix language page: \url{https://docs.verbix.com/Languages/Tangga}
\end{itemize}
The first grammar was written in German by Heinrich Maurer (per the search summary; the work itself was not located).

\section{Religious and Community Resources}
None found.

\section{Digital and Audio Resources}
None found.

\section{Curriculum Development Context}
The language is written in the Roman alphabet. No school materials were found, and the divergence among varieties would need to be resolved with communities before a curriculum is designed.

\end{document}